\section{Introduction}

The full Fourier--cubic quantization of the homogeneous area-preserving
H\'enon map produces exact finite-field trace moments.  Pairing an additive
character with its inverse, then averaging over the maximal real cyclotomic
field, gives rational logarithmic moments \(c_{p,n}\).  Earlier stages of the
construction identify the second moment with a genus-four curve, prove
normal convergence for \(\RePart s>1/3\), and realize the Euler germ by a
sixth-order determinant relative to a normalized faithful semifinite trace.
The first untreated logarithmic term is \(n=3\).

The six chronological steps at \(n=3\) lead naturally to a cubic fourfold
and a \((2,3)\) Fano threefold.  Their middle cohomology supplies the next
square-root gain.

\begin{theorem}[Fano third moment and fourth abscissa]\label{thm:main}
Let \(p>3\) satisfy \(p\equiv1\pmod3\), and let \(\rho\in\Fp^\times\) have
order three.  Outside the finite set of bad reductions of the explicit
complete intersection constructed below, the third descended H\'enon moment
is
\[
 C_{p,3}=-2-\frac{2A_p}{p^2}-\frac{2B_p}{p},
 \qquad
 c_{p,3}=\frac{2C_{p,3}}{p-1},
\]
where \(A_p\) is the primitive-middle Frobenius trace of a Fermat cubic
fourfold and \(B_p\) is the middle Frobenius trace of a smooth \((2,3)\)
threefold.  Their ranks are \(22\) and \(40\), and hence
\[
 |A_p|\le22p^2,
 \qquad |B_p|\le40p^{3/2},
 \qquad
 |c_{p,3}|\le\frac{92+160\sqrt p}{p-1}.
\]
More precisely,
\[
 A_p=20p^2+pa_p,\qquad B_p=pb_p,
 \qquad |a_p|\le2p,\quad |b_p|\le40\sqrt p,
\]
and consequently
\[
 C_{p,3}=-42-2b_p-\frac{2a_p}{p}.
\]
The finitely many omitted local factors are handled by their unitary-block
bound.  Consequently the normalized Euler germ is holomorphic and nonzero
on \(\RePart s>1/4\), and there
\[
 \mathcal G(s)=
 \exp\!\left(-\sum_{n=1}^{7}\frac{\ell_n(s)}n\right)
 \Det_{8,\tau,\mathrm{gr}}(I-X_s).
\]
Order eight is the least fixed integer order forced by positive
\(L^q(\mathcal M,\tau)\)-membership on the whole half-plane.
\end{theorem}

All chronology and all prime clocks are retained before the geometric
reduction.  The theorem is therefore an operator-and-arithmetic advance for
the H\'enon construction, not a claim that the resulting function already
has the Riemann divisor.  Purity is used in the precise form proved by
Deligne \cite[Th\'eor\`eme~1.6]{Deligne1974}; regularized determinant
background is recorded in \cite[Chapter~9]{Simon2005}.  Classical context
for smooth quadric--cubic Fano threefolds appears in
\cite{BlochMurre1979}.  The diagonal-hypersurface calculation uses Weil's
Jacobi-sum formula \cite{Weil1949}; a cohomological formulation is given by
Br\"unjes \cite[Definition~4.4, Proposition~4.5, Theorem~4.6]{Brunjes2003}.
